% Draft sections for the HW4 writeup -- paste into the course .tex template.
% \todo{...} marks values that need results from Torch runs.
\providecommand{\todo}[1]{\textcolor{red}{[TODO: #1]}}

% ---------------- Top of the document ----------------
% GitHub repo: https://github.com/tiwari-adarsh0830/nlp-a4
% Q7 checkpoint (Google Drive): \todo{link}
% AI-use note: \todo{describe how AI tools were used}

% ---------------- Q2a ----------------
\paragraph{Q2a: Transformation.}
I combine two noise sources that occur in real user-written reviews: word choice
(synonym replacement) and typing errors (keyboard typos). The review is split into
words with NLTK's \texttt{word\_tokenize}. Only \emph{eligible} words are changed:
purely alphabetic tokens longer than three characters. This leaves punctuation,
numbers and short function words untouched, including negations such as
``not'' and ``no'' and short sentiment words such as ``bad'', so the label is
preserved. Each eligible word is processed independently:
\begin{itemize}
  \item \textbf{Synonym replacement ($p=0.3$).} I look up the word in WordNet and keep
  only its single most common sense: the synset in which this exact word form has the
  highest SemCor usage count (\texttt{lemma.count()} $>0$). A synonym is drawn uniformly
  from the other single-word lemmas of that synset, and the original capitalisation is
  kept. Restricting to the most common sense avoids wrong-meaning swaps that come from
  pooling every sense (e.g.\ \textit{watch}$\to$\textit{ticker},
  \textit{boring}$\to$\textit{drilling}), and requiring an exact form match avoids
  dropping inflections (e.g.\ \textit{kept}$\to$\textit{maintain}).
  \item \textbf{Typo ($p=0.1$).} One random letter of the word is replaced by a
  neighbouring key on a QWERTY keyboard (e.g.\ \textit{friends}$\to$\textit{friehds}),
  keeping its case.
\end{itemize}
Both steps can apply to the same word. The result is realistic because people
naturally paraphrase with near-synonyms (\textit{movie}/\textit{film},
\textit{boring}/\textit{tedious}, \textit{funny}/\textit{amusing}), and
adjacent-key slips are the most common kind of typing error. A human reader still
recovers the sentiment immediately.

% ---------------- Q3 ----------------
\paragraph{Q3: Data augmentation.}
\todo{Original-train model: accuracy on original test = X\%, on transformed test = Y\%.
Augmented model: original test = X'\%, transformed test = Y'\%.
Did augmentation help on the transformed set? What happened on the original set?
Intuition: training on 5k transformed reviews exposes the model to synonym/typo
variants, so it relies less on exact surface forms. Limitation: it only covers the
specific transformation it was trained on (e.g.\ it will not help with other
distribution shifts) and the 5k examples are noisy copies of reviews already in training.}

% ---------------- Q4 ----------------
\paragraph{Q4: Data statistics.} All lengths are in T5 (\texttt{google-t5/t5-small})
tokens; vocabulary size is the number of distinct token ids used.

\begin{table}[h]
\centering
\begin{tabular}{lcc}
\toprule
Statistics Name & Train & Dev \\
\midrule
Number of examples & 4225 & 466 \\
Mean sentence length & 17.10 & 17.07 \\
Mean SQL query length & 216.37 & 210.05 \\
Vocabulary size (natural language) & 791 & 465 \\
Vocabulary size (SQL) & 555 & 395 \\
\bottomrule
\end{tabular}
\caption{Data statistics before any pre-processing (raw text tokenized with the T5
tokenizer, no \texttt{</s>}). Note that 393 train and 45 dev SQL queries contain
\texttt{<unk>}, because ``\texttt{<}'' is not in T5's vocabulary.}
\end{table}

\begin{table}[h]
\centering
\begin{tabular}{lcc}
\toprule
Statistics Name & Train & Dev \\
\midrule
\multicolumn{3}{l}{\textbf{Model name:} T5 fine-tuned (\texttt{google-t5/t5-small})} \\
Mean sentence length & 18.10 & 18.07 \\
Mean SQL query length & 217.37 & 211.05 \\
Vocabulary size (natural language) & 792 & 466 \\
Vocabulary size (SQL) & 556 & 396 \\
\bottomrule
\end{tabular}
\caption{Data statistics after pre-processing (``\texttt{<}'' replaced by
\texttt{LT}, \texttt{</s>} appended). No \texttt{<unk>} tokens remain. The longest SQL
target is 511 tokens (train) / 503 (dev). Test: 432 examples, mean question length
15.77 tokens.}
\end{table}

% ---------------- Q5 ----------------
\paragraph{Q5: T5 fine-tuning design.}
\begin{table}[h]
\centering
\begin{tabular}{p{3cm}p{10cm}}
\toprule
Design choice & Description \\
\midrule
Data processing &
Natural-language questions are used as-is. In SQL targets, ``\texttt{<}'' is replaced by
the placeholder word \texttt{LT} before tokenization and mapped back after decoding: the
T5 vocabulary has no ``\texttt{<}'', so it would become \texttt{<unk>} and vanish,
making $\sim$10\% of gold queries impossible to generate. \texttt{LT} never occurs in the
data, and every train/dev query round-trips exactly after this mapping. \\
Tokenization &
Encoder and decoder both use the \texttt{google-t5/t5-small} tokenizer (which appends
\texttt{</s>}). The decoder input is \texttt{<extra\_id\_0>} followed by the SQL tokens
minus the last one, and the target is the full SQL sequence ending in \texttt{</s>}.
Batches are padded dynamically with id 0, with an encoder attention mask, and padded
target positions are excluded from the loss. Generation starts from the same
\texttt{<extra\_id\_0>} token (greedy decoding, up to 600 new tokens) and is decoded
with special tokens skipped. \\
Architecture &
Full fine-tuning of the pretrained T5-small: all layers (shared embeddings, all encoder
and decoder blocks, and the tied LM head; 60.5M parameters) are trained. No layers are
frozen. \\
Hyperparameters &
AdamW, learning rate \todo{}, weight decay \todo{}, \todo{cosine} schedule with
\todo{} warm-up epoch(s), batch size \todo{16}, up to \todo{} epochs with early stopping
on dev record F1 (patience \todo{}). \\
\bottomrule
\end{tabular}
\caption{Details of the best-performing T5 model configuration.}
\end{table}

% ---------------- Q6 ----------------
\paragraph{Q6: Results and error analysis.}
\todo{Table 4: dev and test Query EM / record F1 for the final model, plus each variant
tried (e.g.\ learning rate, epochs, with/without the ``<'' fix). Table 5: at least 3
error types with an example and count (e.g.\ wrong/missing table in FROM, wrong
constant such as city or time, missing date/day condition, syntax errors / truncated
query), from \texttt{results/t5\_ft\_experiment\_dev.sql}.}
